\cotefile{136}
\scanpage{80}
\margin{Faisceautisation du topos De Rham}
1. Soit $X$ un topos, et
\[
\mathbb{Z}\xrightarrow{u}\Phi
\]
une immersion de $\mathbb{Z}$ dans un faisceau $\Phi$, tel que

a) $\operatorname{Coker}\Phi$ est $\mathbb{Z}$-plat

b) Pour tout $i>0$, $\Lambda_{\mathbb{Z}}^i\Phi$ et $\Gamma_{\mathbb{Z}}^i\Phi$ sont flasques, et le restent par tensorisation avec tout faisceau abélien.

En particulier, on en conclut \inserted{\deleted{par \unsure{récurrence} et dévissage}} que \deleted{\unsure{si}}

b') Si $i,j\geq 0$ et $i+j>0$, alors $\deleted{\Lambda}\Gamma_{\mathbb{Z}}^i\Phi\otimes\Lambda_{\mathbb{Z}}^j\Phi$ est flasque.

De b') on conclut \deleted{que}\inserted{(par récurrence et dévissage)}, posant
\[
\operatorname{Coker}i=\Psi
\]
\inserted{que} $\alpha$) les $\Gamma_{\mathbb{Z}}^i\Phi\otimes\Lambda_{\mathbb{Z}}^j\Psi$ sont flasques pour $i>0$. \unsure{$\vert$|1}

De $\unsure{\alpha|a|\text{ceci}}$) on conclut\unsure{:}\deleted{\unsure{que}}

$\beta$) Posant $K_u^{\bullet\bullet}=\Gamma_{\mathbb{Z}}^*\Phi\otimes\Lambda_{\mathbb{Z}}^*\Psi$, avec l'op. diff. \deleted{\unsure{$\delta$}} $d$ compatible aux puissances divisées qui, sur $\deleted{\unsure{\Gamma^0\otimes\Lambda^j}}\Psi$ est nul et sur $\unsure{\overline{\Phi}|\Phi}$ est l'hom canonique $\Phi\longrightarrow\Psi$, et définissant
\[
\mathbb{Z}\{T\}\longrightarrow\mathscr{H}^{*0}(K_u^{**})
\]
par $T\longmapsto 1\in\mathbb{Z}=\operatorname{Ker}(\Phi\xrightarrow{d}\Psi)$, et prolongeant par la condition d'être compatible aux p. div, on trouve que $K_u^{\bullet\bullet}$ est une \uline{résolution} de $\mathbb{Z}\{T\}$.

Utilisant encore a), \inserted{et b),} on trouve

$\gamma$) Pour tout faisceau d'anneaux comm. $k$ sur $X$, $K_{u,k}^{\bullet\bullet}\simeq K_u^{\bullet\bullet}\otimes k\simeq\Gamma_k^*\Phi_k\otimes\Lambda_k^*\Psi_k$ est une résolution de $k\{T\}\simeq\mathscr{H}^{\bullet 0}(K_{u,k}^{\bullet\bullet})$, et les $\Gamma_k^i\Phi_k\otimes\Lambda_k^j\Psi_k$ sont flasques pour $i>0$.
% archive sheet number: 79
